\documentclass[11pt]{article}

\usepackage[margin=1in]{geometry}
\usepackage{amsmath,amssymb}
\usepackage{booktabs}
\usepackage{microtype}
\usepackage{hyperref}
\usepackage{xcolor}
\usepackage{enumitem}
\usepackage{array}

\hypersetup{
  colorlinks=true,
  linkcolor=black,
  urlcolor=blue!60!black,
}

\title{AniSync Benchmark: Methodology and Findings}
\author{Tony Dong \and Edward Kim \and Jakob Rees \and Vincent Wang \and Jack Xiong}
\date{April 2026}

\begin{document}
\maketitle

\section{Why benchmark at all}

AniSync is a group anime-discovery app: several users join a room, each
expresses what they want to watch, and the system returns a shared
recommendation list. The product goal is \emph{coherent exploration} ---
surfacing anime the group probably has not seen but will collectively
enjoy --- not retrieval of titles any one user has already rated.

A benchmark cannot directly measure that goal. There is no public dataset
of ``what would this group of four strangers happily watch together if
asked today.'' What we do have is MyAnimeList rating data: per-user
ratings on a 1--10 scale, against an anime catalog that overlaps ours.
The benchmark uses this data to ask a narrower question:

\begin{quote}
Given half of a user's rated anime as visible context, how well does an
algorithm rank a candidate pool against the \emph{other} half of that
user's ratings?
\end{quote}

This is a recall task, not a discovery task. We are explicit about the
mismatch in Section~\ref{sec:limits}. The benchmark's purpose is
threefold: \emph{rule out clearly broken approaches} (sieve), \emph{confirm
that surviving algorithms coherently track individual preference signal}
(coherence floor), and \emph{provide an apples-to-apples comparison between
candidate scoring functions} (ablation surface). The coherence floor
matters because the product arguments for group-level fairness and
diversity only make sense given that the underlying preference-tracking
works at all --- fairness in distributing irrelevant recommendations
is not a feature. The benchmark is not designed to choose the final
algorithm on its own.

\section{The metric: NDCG@K}

\subsection{Definition}

For a single user with a set of relevant items having graded relevances
$r_1, r_2, \dots$, and an algorithm that returns a ranked list of
candidates $c_1, c_2, \dots, c_K$, the Discounted Cumulative Gain at $K$
is
\begin{equation}
\mathrm{DCG@}K \;=\; \sum_{i=1}^{K} \frac{\mathrm{rel}(c_i)}{\log_2(i+1)},
\end{equation}
where $\mathrm{rel}(c_i)$ is the user's relevance score for item $c_i$
(0 if $c_i$ is not in the relevant set). The Ideal DCG is the same sum
computed against the best possible ordering (relevances sorted in
descending order):
\begin{equation}
\mathrm{IDCG@}K \;=\; \sum_{i=1}^{K} \frac{r_{(i)}}{\log_2(i+1)},
\end{equation}
with $r_{(1)} \geq r_{(2)} \geq \dots$ the sorted relevances. NDCG is the
normalised ratio:
\begin{equation}
\mathrm{NDCG@}K \;=\; \frac{\mathrm{DCG@}K}{\mathrm{IDCG@}K} \;\in\; [0, 1].
\end{equation}
A perfect ranking gets 1; a ranking with no relevant items in the top $K$
gets 0.

\subsection{Worked example (binary relevances)}

Suppose a user's relevant set is $\{A, B, C, D\}$, each with relevance 1
(everything else relevance 0), and we evaluate at $K = 5$. The algorithm
returns the ranked list
\[
[\, A,\; X,\; B,\; Y,\; C\, ],
\]
which has relevances $[1, 0, 1, 0, 1]$.

\paragraph{DCG@5.}
\[
\mathrm{DCG@}5 \;=\; \frac{1}{\log_2 2} + \frac{0}{\log_2 3} + \frac{1}{\log_2 4} + \frac{0}{\log_2 5} + \frac{1}{\log_2 6}
\;=\; 1.000 + 0 + 0.500 + 0 + 0.387 \;\approx\; 1.887.
\]

\paragraph{IDCG@5.}
The ideal ranking puts the four relevant items first:
$[1, 1, 1, 1, 0]$.
\[
\mathrm{IDCG@}5 \;=\; \frac{1}{\log_2 2} + \frac{1}{\log_2 3} + \frac{1}{\log_2 4} + \frac{1}{\log_2 5} + 0
\;\approx\; 1.000 + 0.631 + 0.500 + 0.431 + 0 \;=\; 2.562.
\]

\paragraph{NDCG@5.}
\[
\mathrm{NDCG@}5 \;=\; \frac{1.887}{2.562} \;\approx\; 0.737.
\]

The algorithm captured three of four relevant items in the top five, but
because two non-relevant items pushed the third hit down to rank 5, it
loses about a quarter of the ideal score. NDCG punishes mid-rank hits
roughly logarithmically: had the ranking been
$[A, B, X, Y, C]$ instead, DCG@5 would rise to
$1 + 0.631 + 0 + 0 + 0.387 = 2.018$ and NDCG@5 would climb to $0.788$.

\subsection{Why NDCG and not just precision or recall}

Recall@$K$ counts hits without caring about position: any of the rankings
$[A,B,C,X,Y]$, $[A,X,B,Y,C]$, $[X,Y,A,B,C]$ scores identically. For a
recommendation list a user reads top-down, this is wrong. Precision@$K$
suffers the same flaw, and additionally penalises the algorithm for
having more than $K$ correct items available. NDCG's logarithmic
discount aligns better with how users actually consume ranked lists, and
its normalisation produces a $[0, 1]$ score comparable across users with
different numbers of relevant items.

\section{Setup}

\paragraph{User pool.}
2{,}000 MyAnimeList users sampled from rating-export data, filtered to
those with at least 15 rated titles. Ratings are integers 1--10 (we drop
the 0 ``not rated yet'' value).

\paragraph{Catalog.}
The full AniSync catalog (anime-offline-database, $\sim$33{,}000 entries
post-filter), preprocessed into 384-dimensional cosine-normalised
embeddings via \texttt{sentence-transformers/all-MiniLM-L6-v2} on a
constructed \texttt{text\_blob} per item (title, synonyms, tags, type,
year, season, studios, producers, score).

\paragraph{Visible/hidden split.}
For each user we permute their ratings deterministically by
$(\text{username}, \text{profile\_seed})$ and split 50/50. The
\emph{visible} half is what the algorithm sees as input. The
\emph{hidden} half is the ground truth used to evaluate the output.

\paragraph{Groups.}
We sample 400 synthetic groups of 4 users uniformly at random with a
fixed \texttt{group\_seed}. Group composition is independent of taste --- a
deliberate stress test for the group-aggregation logic.

\paragraph{Evaluation.}
For each group we run the algorithm on the four users' visible ratings,
take the top 200 ranked catalog items, and compute per-user NDCG@200
against that user's relevance set (Section~\ref{sec:proxy}). The group
score is the mean across the four users; we also track the
worst-user-in-group NDCG. We report the mean and standard deviation of
both quantities across the 400 groups.

\section{Proxy relevance set}
\label{sec:proxy}

\subsection{The temporal gap}

The ratings dataset is from 2018; the catalog runs to 2026. A user who
loved \emph{Steins;Gate} should plausibly enjoy a thematically similar
2022 release, but they cannot have rated it. Scoring against the literal
hidden set under-credits any algorithm that surfaces post-2018 anime ---
exactly the most useful behaviour for a discovery product.

\subsection{Construction}

For each user we expand the literal hidden set into a \emph{proxy
relevance set} as follows:
\begin{enumerate}[leftmargin=*]
\item Take the top $n_{\mathrm{top}} = 10$ hidden items by score, keeping
only those with score $> 5$ (positive signal).
\item For each seed item $s$ with score $\sigma_s$, retrieve its
$n_{\mathrm{neighbours}} = 10$ nearest catalog entries by cosine
distance.
\item Assign each neighbour $c$ a relevance
$\mathrm{rel}(c) = \max(\sigma_s - 5, 0)$, taking the maximum across
seeds when multiple seeds nominate the same neighbour.
\end{enumerate}
This yields up to 100 graded-relevance items per user. The relevance
scale ranges from 0 (not relevant or score $\leq 5$) to 5 (a
neighbour of a 10/10).

\subsection{What this metric does and does not measure}

The proxy set rewards algorithms that surface catalog items
\emph{semantically close} to a user's known high-score items. It does
not reward genuine novelty: an algorithm that recommends anime
unrelated to anything in the user's history will score zero, even if
the user would in fact enjoy them. It also does not measure fairness;
NDCG averaged across users in a group cannot distinguish ``everyone
gets one decent recommendation'' from ``three people get great results
and one gets nothing.''

We return to this in Section~\ref{sec:limits}.

\section{Algorithms compared}
\label{sec:algos}

Six algorithms are evaluated. Each is described in one paragraph below
with an ``isolates'' tag identifying what it tests.

\paragraph{\texttt{centroid}.}
Compute a single mean embedding from all liked items across all four
users. Retrieve the nearest catalog items by cosine to that centroid;
return them in distance order. \emph{Isolates: the value of treating the
group as a single average user.}

\paragraph{\texttt{groupmatch\_raw}.}
For each user, build a query embedding as the mean of their liked-item
embeddings (rating $\geq 7$, falling back to all visible if none
qualify). Retrieve each user's top-100 candidates via pgvector and union
the four pools. Score each candidate by the \emph{mean} cosine
similarity across the four user-query embeddings. Return the top 200 by
mean score. \emph{Isolates: per-user retrieval plus simple averaging,
without clustering or fairness machinery.}

\paragraph{\texttt{groupmatch\_clustered}.}
Identical to \texttt{groupmatch\_raw}, but after scoring we run
silhouette-guided $k$-means on the candidate embeddings, promote the
top-$N$ items per cluster to the front of the ranked list, and globally
re-sort the promoted set by mean GroupMatch score. Items not promoted
follow, also in score order. \emph{Isolates: the marginal value of
cluster-diverse re-rank on top of mean-similarity scoring.}

\paragraph{\texttt{groupfit}.}
For each candidate $c_i$, score
\[
\mathrm{score}(c_i) \;=\; \min_{u} \Big[ \max_{j} (e_i \cdot \ell_{u,j})
\;-\; \lambda\, (e_i \cdot \bar n_u)
\;+\; \beta\, (e_i \cdot t_u) \Big],
\]
where $\ell_{u,j}$ are user $u$'s liked-item embeddings, $\bar n_u$ is
the mean of their disliked-item embeddings (negative centroid), and $t_u$
is a query embedding from their visible profile. The min over users
encodes a fairness constraint: a candidate scores well only if every user
in the group has at least one taste vector matched. \emph{Isolates: min-aggregation
fairness combined with an additive negative-penalty term and a text-alignment
term.} The benchmark grids over $(\lambda, \beta)$.

\paragraph{\texttt{groupfit\_pos\_text}.}
A simplification of \texttt{groupfit} that drops the negative term and
blends the positive and text terms with a single mixing parameter:
\[
\mathrm{score}(c_i) \;=\; (1 - \alpha)\, \min_{u} \max_{j} (e_i \cdot \ell_{u,j})
\;+\; \alpha\, \frac{1}{|U|} \sum_{u} (e_i \cdot t_u).
\]
\emph{Isolates: text alignment as a clean blend against positive-only
fairness, with no negative penalty.} The benchmark sweeps $\alpha$.

\paragraph{\texttt{groupmatch\_raw\_llm}.}
Identical to \texttt{groupmatch\_raw}, but the per-user query embedding
is computed from an LLM-generated natural-language description of the
user's taste (built offline from their visible ratings) rather than from
the mean of their liked-item embeddings. \emph{Isolates: whether
LLM-mediated taste extraction beats raw embedding averaging.}

\paragraph{The \texttt{\_msmarco} variants.}
\texttt{groupfit\_msmarco} and \texttt{groupfit\_pos\_text\_msmarco} are
the same algorithms with the text term computed against an asymmetric
model (\texttt{msmarco-MiniLM-L6-cos-v5}) populated into a separate
\texttt{embedding\_msmarco} column. The hypothesis was that an
asymmetric query$\to$document model would outperform the symmetric
\texttt{all-MiniLM-L6-v2} on the natural-language preference path. The
benchmark refutes it (Section~\ref{sec:rejected}).

\section{Results}
\label{sec:results}

All runs: 400 groups of 4 users, visible ratio 0.5, $K = 200$, profile
seed and group seed fixed for reproducibility.

\subsection{Algorithm comparison}

\begin{center}
\begin{tabular}{lccc}
\toprule
Algorithm & mean NDCG & std & mean worst-user NDCG \\
\midrule
\texttt{groupmatch\_clustered}      & 0.1065 & 0.0276 & 0.0538 \\
\texttt{groupmatch\_raw}            & 0.1058 & 0.0272 & 0.0528 \\
\texttt{groupfit\_pos\_text}        & 0.1024 & 0.0286 & 0.0509 \\
\texttt{groupfit}                   & 0.1006 & 0.0273 & 0.0492 \\
\texttt{groupfit\_msmarco}          & 0.0944 & 0.0260 & 0.0481 \\
\texttt{groupfit\_pos\_text\_msmarco} & 0.0941 & 0.0247 & 0.0482 \\
\texttt{centroid}                   & 0.0798 & 0.0235 & 0.0342 \\
\texttt{groupmatch\_raw\_llm}       & 0.0470 & 0.0170 & 0.0188 \\
\bottomrule
\end{tabular}
\end{center}

\paragraph{Reading the absolute numbers.}
NDCG of $\sim 0.10$ looks low at face value, but the right comparison
is to the random baseline. With $\sim 50$ relevant items per user in a
$\sim 33{,}000$-item catalog, random ranking scores near zero on this
metric. The fact that the surviving algorithms cluster in the
$0.10$ range --- and that \texttt{groupmatch\_raw\_llm} at $0.047$ is
visibly distinguishable as broken --- means the metric is sensitive
enough to discriminate algorithms that track preference from those
that do not. Differences between close performers (e.g.\ the $\sim 4\%$
gap between \texttt{groupmatch\_clustered} and \texttt{groupfit\_pos\_text})
are within the per-group standard deviation and should be interpreted
accordingly.

\subsection{GroupFit ablation: $\lambda \times \beta$}

Mean NDCG across the negative-penalty weight $\lambda$ (rows) and the
text-alignment weight $\beta$ (columns), with all other parameters
fixed.

\begin{center}
\begin{tabular}{lccccccc}
\toprule
$\lambda \backslash \beta$ & 0.0 & 0.1 & 0.2 & 0.3 & 0.5 & 0.7 & 1.0 \\
\midrule
0.0 & 0.0974 & 0.0969 & 0.0964 & 0.0957 & 0.0941 & 0.0927 & 0.0909 \\
0.2 & 0.0990 & 0.0988 & 0.0983 & 0.0975 & 0.0961 & 0.0947 & 0.0929 \\
0.5 & 0.1000 & 0.0998 & 0.0995 & 0.0990 & 0.0981 & 0.0967 & 0.0947 \\
0.8 & 0.1002 & 0.1000 & 0.0997 & 0.0993 & 0.0984 & 0.0974 & --     \\
1.0 & 0.1001 & 0.0999 & 0.0996 & 0.0992 & 0.0983 & 0.0977 & --     \\
\bottomrule
\end{tabular}
\end{center}

The grid is monotonically decreasing left-to-right (in $\beta$) and
roughly increasing top-to-bottom (in $\lambda$).

\subsection{GroupFit pos+text ablation: $\alpha$}

\begin{center}
\begin{tabular}{lccc}
\toprule
$\alpha$ & mean NDCG & std & mean worst-user NDCG \\
\midrule
0.0 & 0.0979 & 0.0276 & 0.0490 \\
0.1 & 0.0994 & 0.0279 & 0.0495 \\
0.2 & 0.1004 & 0.0285 & 0.0499 \\
\textbf{0.3} & \textbf{0.1009} & 0.0286 & \textbf{0.0501} \\
0.4 & 0.1008 & 0.0287 & 0.0498 \\
0.5 & 0.1006 & 0.0290 & 0.0496 \\
0.6 & 0.0997 & 0.0285 & 0.0491 \\
0.7 & 0.0990 & 0.0286 & 0.0488 \\
0.8 & 0.0981 & 0.0282 & 0.0484 \\
0.9 & 0.0972 & 0.0278 & 0.0479 \\
1.0 & 0.0963 & 0.0273 & 0.0474 \\
\bottomrule
\end{tabular}
\end{center}

The optimum sits at $\alpha \approx 0.3$. Pure liked-item retrieval
($\alpha = 0$) and pure text alignment ($\alpha = 1$) both underperform
the blend.

\section{Three rejected hypotheses}
\label{sec:rejected}

The benchmark cleanly disposes of three design candidates that we had
considered serious going in.

\paragraph{Negative-similarity penalty ($\lambda > 0$).}
Penalising similarity to a centroid of the user's disliked items was
intuitive --- ``don't recommend things they have already rejected.''
The $\lambda \times \beta$ grid shows it is monotonically harmful: at
every fixed $\beta$, raising $\lambda$ from 0 reduces NDCG. The mean
of disliked items is too diffuse a target; subtracting from candidate
scores along that direction discards positive signal as collateral.
\textbf{Decision: $\lambda = 0$ in production.}

\paragraph{Centroid retrieval.}
Treating the group as a single user (mean of liked items across the
group, retrieve nearest neighbours) loses 20--25\% NDCG relative to per-user
retrieval. Averaging four taste vectors before retrieval collapses the
group's structure; per-user retrieval followed by group-aware scoring
preserves it.
\textbf{Decision: per-user retrieval, never group-mean retrieval.}

\paragraph{Asymmetric msmarco embeddings.}
\texttt{msmarco-MiniLM-L6-cos-v5} is fine-tuned on
query$\to$passage pairs and was hypothesised to better match
preference-text$\to$anime-metadata than the symmetric
\texttt{all-MiniLM-L6-v2}. In practice, both \texttt{groupfit\_msmarco}
(0.0944) and \texttt{groupfit\_pos\_text\_msmarco} (0.0941) score about
6\% below their non-msmarco counterparts (0.1006 and 0.1024). The
proxy relevance set is constructed via cosine neighbours under
\texttt{all-MiniLM-L6-v2}; an algorithm that scores items with a
\emph{different} embedding cannot align as well with the metric's
notion of ``relevant.'' Whatever asymmetric retrieval gains exist are
not detectable by this benchmark, and we have no other source of
evidence for them.
\textbf{Decision: a single embedding model in production
(\texttt{all-MiniLM-L6-v2}). The msmarco column and embedding code stay
in the repo as the executed ablation but are not invoked by the
production path.}

\section{What we keep}

\paragraph{Text alignment is a real, modest signal.}
Within \texttt{groupfit\_pos\_text}, $\alpha = 0$ scores 0.0979 and
$\alpha = 0.3$ scores 0.1009 --- a 3\% relative gain. The text term
brings information the liked-item embeddings cannot: it lets a user
type ``something slow-paced and melancholic'' and have that route their
recommendations regardless of what is in their history.

\paragraph{Cluster-diverse re-rank is worth the complexity.}
\texttt{groupmatch\_clustered} (0.1065) edges out
\texttt{groupmatch\_raw} (0.1058). The gain is small in NDCG terms but
consistent, and the product reason for clustering --- giving a session
visible diversity rather than a uniform top-200 by score --- is itself
a goal the benchmark cannot directly reward.

\paragraph{Per-user retrieval beats centroid retrieval.}
Confirmed by the centroid result above. Production retrieval is one
pgvector query per user, candidate pools unioned.

\section{Limitations: what the benchmark cannot adjudicate}
\label{sec:limits}

\paragraph{The metric is constructed using the same machinery as
GroupMatch retrieval.}
The proxy relevance set is built by taking nearest embedding neighbours
of each user's high-rated hidden items. \texttt{groupmatch\_raw}
retrieves candidates by taking nearest embedding neighbours of each
user's high-rated \emph{visible} items and scores them by mean cosine.
The metric is, by construction, a target that GroupMatch is shaped
to hit. The 6\% gap between GroupMatch and GroupFit on this benchmark
is not strong evidence that GroupMatch is the better algorithm in
production --- it is partly evidence that the metric and the algorithm
share an inductive bias.

\paragraph{NDCG measures average satisfaction, not fairness in the room.}
GroupFit's $\min_u$ aggregation deliberately sacrifices group-mean
quality to ensure every individual user has at least one taste direction
hit. NDCG, averaged over users, cannot capture this property. A group
where one user is entirely unrepresented and three are well served
will outscore a group where all four are moderately served. Worst-user
NDCG (which we report) is closer to a fairness signal but still does
not capture the qualitative experience of being the one person in the
room with no candidate they connect to.

\paragraph{The benchmark rewards recall, not discovery.}
Every algorithm in this benchmark is graded on how well it surfaces
items semantically close to what a user already liked in 2018. There
is a trivial way to maximise this metric: recommend the most popular,
highest-rated anime in the catalog, since those items are
disproportionately likely to be in any user's neighbourhood. An
algorithm that did exactly that would be a worse product. AniSync's
goal is the opposite of recall: surface anime the group has not
already encountered. The benchmark is structurally hostile to that
goal.

\paragraph{Ratings are eight years stale.}
The proxy expansion partially compensates for this, but a user's
2018 preferences and their 2026 preferences are not the same data,
and the benchmark cannot model the difference.

\section{From benchmark to product}
\label{sec:design}

The benchmark told us three things to reject (negative penalty,
centroid retrieval, msmarco embeddings) and one signal that helps (text
alignment, around $\alpha \approx 0.3$). It did not tell us whether
GroupMatch or GroupFit is the right algorithm for the product, because
the metric is incapable of measuring the property that distinguishes
them: \emph{room-level fairness} and \emph{willingness to recommend
unfamiliar anime instead of crowd-favourites}.

We choose \texttt{groupfit\_pos\_text} despite the benchmark numbers
because:

\begin{itemize}[leftmargin=*]
\item \textbf{Fairness in the room.} A group session where one
participant is completely unrepresented is a product failure. Min-over-users
encodes this constraint directly. Mean-over-users does not.

\item \textbf{Exploration over recall.} Per-user max-similarity to
liked items is more discriminating than mean similarity to a group
centroid: it surfaces specific items that hit \emph{some} user's
taste vector hard, rather than generically popular items that
moderately match everyone. This is the right inductive bias for
discovery.

\item \textbf{Text as a first-class signal.} A user's typed preference
(``I want something experimental and short'') is information the
embedding of their watch history cannot represent. The benchmark's
$\alpha$ ablation confirms that even a small text weight measurably
improves alignment.
\end{itemize}

\subsection*{Production configuration}

\begin{itemize}[leftmargin=*]
\item \textbf{Algorithm:} \texttt{groupfit\_pos\_text} with
$\alpha = 0.30$, $\lambda = 0$.
\item \textbf{Embedding model:} \texttt{all-MiniLM-L6-v2}, 384 dim,
cosine-normalised, single column on \texttt{catalog\_items}.
\item \textbf{Retrieval:} per-user pgvector query against the mean of
each user's liked-item embeddings (score $\geq 7$, falling back to all
visible if none qualify), top 100 per user, candidate pools unioned.
\item \textbf{Diversification:} silhouette-guided $k$-means on the
candidate pool, top-$N$ per cluster promoted to the front of the
ranked list, globally re-sorted by GroupFit score.
\item \textbf{Negative ratings:} not used in scoring. Empirically
harmful, theoretically suspect, removed.
\end{itemize}

The benchmark is the floor, not the ceiling: it confirms our scoring
function is in the right family and our preprocessing choices are
not broken. The product case for the chosen design rests on properties
the benchmark cannot see.

\end{document}